\section{Transformer Comparison}
\label{sec:transformer-comparison}

Traditional transformer systems
rely on memoryless self-attention:
\[
\text{Attention}(Q,K,V) = 
\text{softmax}\left(
\frac{QK^T}{\sqrt{d_k}}
\right)V
\]
This operational matrix assumes
a flat Euclidean embedding space
where individual token vectors
undergo linear projections.

Because the softmax normalization
lacks path-dependent memory,
these models exhibit high
susceptibility to semantic mimicry
and context drift. The model
completes token patterns without
tracking geometric deformation.

In contrast, the NDH operator spine
replaces linear attention weights
with continuous, path-ordered
holonomy transport steps:
\[
\mathbf{z}_{\text{return}} 
= \mathcal{P} \exp\left(
\oint_{\gamma} \Gamma^\mu_{\nu\lambda} 
\, dx^\lambda \right) \mathbf{z}_0
\]
By routing dynamic text state
updates along non-Euclidean curves
bound by the 9D radial Laplacian,
the NDH architecture prevents
the metric collapse anomalies
inherent in commercial scaling.
